\documentclass[border=10pt]{standalone}
\usepackage{circuitikz}
\usetikzlibrary{calc}
\begin{document}
\begin{circuitikz}[scale=1.2, transform shape]
  % --- transistors ---
  \node[pmos, yscale=-1] (P) at (0,2.1) {};   % flipped: source on top
  \node[nmos]            (N) at (0,-0.4) {};

  % --- supply rails ---
  \draw (P.S) -- (0,3.3) node[vcc]{$V_{DD}$};
  \draw (N.S) -- (0,-1.6) node[ground]{};

  % --- output column (drain to drain) ---
  \draw (P.D) -- (N.D);
  \coordinate (mid) at ($(P.D)!0.5!(N.D)$);
  \draw (mid) -- ++(2.2,0) coordinate (vo);   % horizontal tap (no slant)
  \node[ocirc] at (vo) {};
  \node[right] at (vo) {$V_{out}$};

  % --- input (gates tied) ---
  \draw (P.G) -- (-1.3, 2.1);
  \draw (N.G) -- (-1.3,-0.4);
  \draw (-1.3,2.1) -- (-1.3,-0.4);
  \coordinate (inmid) at (-1.3,0.85);
  \draw (inmid) -- (-2.4,0.85);
  \node[ocirc] at (-2.4,0.85) {};
  \node[left] at (-2.5,0.85) {$V_{in}$};

  % labels
  \node[right] at (P.D) {\scriptsize PMOS};
  \node[right] at (N.D) {\scriptsize NMOS};
\end{circuitikz}
\end{document}
